\section{Results}
\label{sec:results}

Results cover the twelve Project-Trained Configurations and the 91 admitted
model-variant groups of Figure~\ref{fig:deployment-admission}: 22 same-lane
controls and 69 transformed groups, measured over 273 executions.  All
quality values are the estimator of Eq.~\eqref{eq:quality-estimator} on the
fixed adaptive LJSpeech evaluation set, and every claim concerns these executed
states.  Except for static ONNX INT8
(Section~\ref{sec:deployment-methods}), the three executions per group
repeat inference from one frozen artifact, so their spread reflects
execution variability rather than statistical uncertainty.

\subsection{Baseline Quality, Efficiency, and Footprint}
\label{sec:baseline-results}

Addressing RQ1, Table~\ref{tab:baseline-surface} reports the complete
twelve-configuration outcome surface, including the floor-level negative
results, and
Figure~\ref{fig:baseline-landscape} places every configuration on the
quality-speed plane.  True-length PESQ ranked RNDVoC (3.895), RFWave
(3.785), and Vocos (3.614) as the three strongest states, and STOI produced
the same leading order; seven configurations reached true-length PESQ of at
least 3.30.  UTMOS disagreed with the intrusive proxies: HiFi-GAN V1 was
highest at 4.203, followed by Vocos (4.089) and BigVGAN-base (4.062); no
configuration was uniformly best across the three proxies.

\begin{table}[t]
  \centering
  \footnotesize
  \renewcommand{\arraystretch}{1.02}
  \setlength{\tabcolsep}{4pt}
  \InputIfFileExists{sections/results/scripts/generated/baseline_table.tex}{}{%
    \begin{tabular*}{\linewidth}{@{\extracolsep{\fill}}lrrrrrrr@{}}
  \toprule
  Configuration & Params & Size & PESQ & STOI & UTMOS & RTF B200 & RTF CPU$^{\mathsection}$ \\
  \midrule
  APNet2 & 31.43 & 119.88 & 2.515 & 0.930 & 3.091 & 0.0095 & 0.0191 \\
  BigVGAN-base & 14.03 & 53.51 & 3.320 & 0.974 & 4.062 & 0.1573 & 0.6985 \\
  FreeV & 18.22 & 69.66 & 3.386 & 0.969 & 3.789 & 0.0077 & 0.0157 \\
  HiFi-GAN V1 & 13.94 & 53.16 & 3.305 & 0.974 & 4.203 & 0.1197 & 0.1836 \\
  HiFi-GAN V2$^{\ddagger}$ & 0.93 & 3.54 & 2.400 & 0.942 & 3.907 & 0.1140 & 0.0507 \\
  HiFi-GAN V3$^{\ddagger}$ & 1.46 & 5.59 & 2.550 & 0.949 & 3.851 & 0.0580 & 0.0366 \\
  LPCNet & 1.23 & 4.70 & 1.052 & 0.527 & 1.235 & 11.36$^{*}$ & 9.2015 \\
  MelGAN & 4.27 & 16.27 & 1.043 & 0.418 & 1.736 & 0.0742 & 0.0456 \\
  RFWave$^{\dagger}$ & 18.14 & 69.21 & 3.785 & 0.981 & 3.514 & 0.0148 & 0.7491 \\
  RNDVoC$^{\dagger}$$^{\ddagger}$ & 3.57 & 14.94 & 3.895 & 0.987 & 4.007 & 0.0559 & 0.0838 \\
  Vocos$^{\dagger}$ & 13.53 & 51.62 & 3.614 & 0.979 & 4.089 & 0.0047 & 0.0098 \\
  VocosFormer & 13.78 & 52.57 & 3.512 & 0.976 & 3.901 & 0.0122 & 0.0105 \\
  \bottomrule
\end{tabular*}
}
  \caption{Baseline outcome surface for all twelve configurations: quality
  proxies (higher is better), parameters (millions), serialized size (MB),
  and warm real-time factors (lower is better).
  $^{\mathsection}$retained comparative CPU record
  (Section~\ref{sec:evaluation-protocols}), not a deployment-phase value;
  $^{\dagger}$B200 speed-quality frontier;
  $^{\ddagger}$parameter- and size-quality frontier;
  $^{*}$LPCNet timing (27 utterances, one execution) is excluded from
  strict frontiers.}
  \label{tab:baseline-surface}
\end{table}

\begin{figure}[t]
  \centering
  \includegraphics[width=\linewidth]{baseline_landscape}
  \caption{Twelve frozen baselines on the true-length PESQ versus B200
  real-time-factor plane (Eq.~\eqref{eq:rtf}).  The legend identifies each
  configuration, ordered by PESQ; marker shape encodes the generation
  family (squares: frame-level Fourier; circles: time-domain GAN;
  triangle: flow matching; diamond: hybrid attention; inverted triangle:
  autoregressive).  Dark-edged markers are frontier members, connected by
  the dashed nondominance chain; the dotted line marks real-time
  synthesis.  LPCNet's single-execution timing is shown but excluded from
  the strict frontier.  Each point is one row of the retained comparative
  record.}
  \label{fig:baseline-landscape}
\end{figure}

No checkpoint minimized all deployment costs.  Vocos had the lowest strict
B200 RTF (0.00473), RNDVoC attained the highest PESQ at 11.8 times that
RTF, and RFWave occupied the intermediate point (3.785 at 0.0148); the
strict speed frontier comprised Vocos, RFWave, and RNDVoC.  HiFi-GAN V2 was the smallest baseline (0.93 million
parameters, 3.54~MB) at PESQ 2.400; the parameter and size frontiers
comprised HiFi-GAN V2, HiFi-GAN V3, and RNDVoC.  Two configurations
produced floor-level negative results: MelGAN (PESQ 1.043 on both lanes)
is consistent with budget-limited non-convergence, and LPCNet (1.052)
lost free-running stability at 5.8\% of its reference recipe while its
PyTorch timing omits the author's optimized sparse C path; both remain
full members of the twelve-model cohort.

The matched VocosFormer state did not improve on its Vocos control: it added
247{,}040 parameters ($+1.83\%$), reduced true-length PESQ by 0.101 and STOI
by 0.0035, and increased B200 RTF $2.58\times$.  Under padded batch-16
evaluation the same checkpoint measured PESQ 2.480 against 3.512 at true
length, a padding-sensitivity failure consistent with contamination through
unmasked global attention rather than intrinsic synthesis collapse; this
one-trajectory comparison is a configuration-level observation, not an
average attention effect.

\subsection{Deployment Transformation Effects}
\label{sec:intervention-results}

Addressing RQ2, the decision rule of Eq.~\eqref{eq:decision-rule} labeled
the 69 transformed groups 19 beneficial, 18 inconclusive, and 32 harmful.
Table~\ref{tab:family-effects} reports the seven intervention families and
their fourteen variants with per-variant label composition, median
same-lane effects, and $\Delta$PESQ ranges; the 29-row exclusion
register documents the 88 attempts producing no admitted group.  FP16
weight casting was the most consistently favorable family: eight of
ten groups were beneficial at half serialized size, with RFWave and
HiFi-GAN V3 reaching $1.955\times$ and $1.874\times$ speedups at preserved
PESQ.

\begin{table}[t]
  \centering
  \footnotesize
  \renewcommand{\arraystretch}{1.02}
  \InputIfFileExists{sections/results/scripts/generated/variant_table.tex}{}{%
    \begin{tabular*}{\linewidth}{@{\extracolsep{\fill}}lrcrcrr@{}}
  \toprule
  Variant & $n$ & B/I/H & Med.\ $\Delta$PESQ & $\Delta$PESQ range & Med.\ $S_a$ & Med.\ $Z_a$ \\
  \midrule
  \multicolumn{7}{@{}l}{\emph{Compilation}} \\
  \quad Default mode & 11 & 4/5/2 & $-0.0000$ & $[-0.000, +0.000]$ & 1.07 & 1.00 \\
  \quad Reduce-overhead mode & 4 & 1/0/3 & $-0.0000$ & $[-0.000, +0.000]$ & 0.40 & 1.00 \\
  \addlinespace[2pt]
  \multicolumn{7}{@{}l}{\emph{FP16}} \\
  \quad Weight cast & 10 & 8/2/0 & $-0.0002$ & $[-0.002, +0.001]$ & 1.08 & 0.50 \\
  \addlinespace[2pt]
  \multicolumn{7}{@{}l}{\emph{INT8}} \\
  \quad Dynamic & 5 & 2/1/2 & $-0.0469$ & $[-0.648, +0.004]$ & 1.30 & 0.40 \\
  \quad Weight-only & 4 & 2/2/0 & $-0.0010$ & $[-0.003, +0.002]$ & 0.92 & 0.41 \\
  \addlinespace[2pt]
  \multicolumn{7}{@{}l}{\emph{ONNX}} \\
  \quad FP32 export & 4 & 1/2/1 & $-0.0000$ & $[-0.000, +0.000]$ & 0.90 & 1.00 \\
  \quad INT8 static & 4 & 0/1/3 & $-0.7701$ & $[-1.329, -0.004]$ & 0.82 & 0.26 \\
  \addlinespace[2pt]
  \multicolumn{7}{@{}l}{\emph{Pruning}} \\
  \quad 30\% sparsity & 4 & 0/0/4 & $-1.0259$ & $[-2.740, -0.148]$ & 0.95 & 1.00 \\
  \quad 50\% sparsity & 11 & 0/1/10 & $-1.0632$ & $[-2.719, +0.000]$ & 1.00 & 1.00 \\
  \quad 70\% sparsity & 4 & 0/0/4 & $-2.2722$ & $[-2.712, -2.014]$ & 1.03 & 1.00 \\
  \addlinespace[2pt]
  \multicolumn{7}{@{}l}{\emph{Recovery}} \\
  \quad 50\% + recovery & 5 & 0/4/1 & $-0.0070$ & $[-0.475, +0.017]$ & 1.11 & 1.00 \\
  \addlinespace[2pt]
  \multicolumn{7}{@{}l}{\emph{ODE steps}} \\
  \quad 8 steps & 1 & 1/0/0 & $-0.0456$ & n/a & 1.20 & 1.00 \\
  \quad 4 steps & 1 & 0/0/1 & $-0.3689$ & n/a & 2.20 & 1.00 \\
  \quad 2 steps & 1 & 0/0/1 & $-1.6363$ & n/a & 4.51 & 1.00 \\
  \bottomrule
\end{tabular*}
}
  \caption{Deployment intervention outcomes over the 69 admitted groups,
  by family and variant: admitted count, beneficial/inconclusive/harmful
  composition under Eq.~\eqref{eq:decision-rule}, median same-lane
  $\Delta$PESQ, speedup $S_a$, and size ratio $Z_a$
  (Eq.~\eqref{eq:deployment-effects}), and the $\Delta$PESQ range across
  the variant's admitted groups; single-group variants report their one
  contrast.}
  \label{tab:family-effects}
\end{table}

Compilation was mode specific: default-mode compilation preserved PESQ in
all eleven groups at speedups from $0.14\times$ to $1.403\times$
(HiFi-GAN V3), whereas the reduce-overhead mode was harmful for three of
four groups (speedups 0.34 to 0.41) and beneficial only for RFWave
($1.361\times$).  ONNX FP32
accelerated HiFi-GAN V2 by $2.071\times$ at $\Delta$PESQ $-0.00013$
(fixed-set 95\% interval $[-0.00036, 0.00010]$) yet split 1/2/1 (B/I/H)
overall.  Dynamic INT8 diverged: beneficial for APNet2 and
Vocos, harmful for FreeV ($-0.123$) and RFWave ($-0.648$).  Weight-only
INT8 preserved PESQ in all four groups but slowed FreeV ($0.851\times$) and
VocosFormer ($0.617\times$).

Static ONNX INT8 was harmful in three of four groups with median PESQ loss
0.770 at median size ratio 0.265.  Because each execution recalibrates on
its own shuffled loader batches, its per-execution PESQ disperses where all
other variants repeat exactly: HiFi-GAN V1 measured 2.003, 2.053, and
1.910 against its execution-invariant ONNX FP32 value of 3.317.  Dense magnitude
masks were harmful in 4/4, 10/11, and 4/4 groups at 30/50/70\% sparsity
(median $\Delta$PESQ $-1.03$, $-1.06$, and $-2.27$); the single non-harmful
group is floor-censored MelGAN.  Recovery returned four of five admitted
50\% groups to within 0.05 PESQ of their controls (BigVGAN-base moved
from $-1.149$ to $+0.017$) but left Vocos harmful ($-0.475$); lacking
the dense-continuation control, the recovery arm supports no causal claim.  RFWave
solver reduction traded monotonically: speedups of 1.201, 2.200, and 4.509
for PESQ losses of 0.046, 0.369, and 1.636 at 8, 4, and 2 steps.

Deployment-phase RTFs follow Eq.~\eqref{eq:rtf} with $J_a{=}3$, not
comparable with Table~\ref{tab:baseline-surface}'s single-pass B200 or
comparative CPU baselines (Section~\ref{sec:evaluation-protocols}).  Across the 273 executions, p50 spanned
1.21 to 3{,}909.35~ms, p95 1.96 to 6{,}138.02~ms, process RSS 6{,}754 to
16{,}991~MiB, and the 24 ONNX session builds spanned 0.0116 to 0.1327~s.  The
deployment costs need not move together: the HiFi-GAN V2 ONNX FP32
artifact halves median synthesis time on CPU while raising process RSS
above its control.  Of the 69 utterance-paired PESQ contrasts, 34 fixed-set intervals were
entirely negative, two entirely positive, and 33 crossed zero.  Despite
non-normality (113 contrasts with absolute skewness above one, 181 with
excess kurtosis above three), bootstrap and Student $t$ signs agreed on
482 of 483 contrasts, an internal agreement diagnostic, not external
validity.

\subsection{Sensitivity of the Decision Taxonomy and Constraint-Aware
Selection}
\label{sec:selection-results}

Across the $3^5=243$ threshold profiles of
Section~\ref{sec:evaluation-protocols}, the beneficial,
harmful, and inconclusive totals ranged over 12 to 24, 29 to 34, and 12 to
27; median and minimum agreement with the default taxonomy were 91.3\%
and 82.6\%.  Stability stratifies by rule structure: 27 of the 28 rows
governed by the dense-mask and MelGAN overrides never changed label
(96.4\%), whereas 25 of the 41 threshold-governed rows never changed
(61.0\%), so the aggregate stability (52 of 69 rows unchanged) partly
reflects deterministic rule structure rather than threshold insensitivity.

Frontier membership depended on the proxy: STOI reproduced the
PESQ frontiers of Section~\ref{sec:baseline-results}, while UTMOS yielded
Vocos then HiFi-GAN V1 against B200 RTF and, against parameters, HiFi-GAN
V2, RNDVoC, Vocos, and HiFi-GAN V1, so proxy choice alters which states appear
efficient.

\begin{table}[t]
  \centering
  \footnotesize
  \renewcommand{\arraystretch}{1.02}
  \InputIfFileExists{sections/results/scripts/generated/selector_table.tex}{}{%
    \begin{tabular*}{\linewidth}{@{\extracolsep{\fill}}llrlrr@{}}
  \toprule
  Constraints $(q_{\min}, P_{\max}, D_{\max})$ & Lane & Feasible & Selected artifact $a_h^\star$ & RTF & PESQ \\
  \midrule
  (3.30, 20M, 70\,MB) & B200 & 28 & FP16-cast Vocos & 0.0003 & 3.611 \\
  (3.30, 20M, 70\,MB) & CPU & 8 & Dynamic-INT8 Vocos & 0.0055 & 3.563 \\
  (3.60, 15M, 55\,MB) & B200 & 10 & FP16-cast Vocos & 0.0003 & 3.611 \\
  (3.60, 15M, 55\,MB) & CPU & 2 & Untransformed Vocos & 0.0067 & 3.610 \\
  \bottomrule
\end{tabular*}
}
  \caption{Constraint-aware selection (Eq.~\eqref{eq:constrained-selection})
  over all admitted artifacts per lane, with same-lane controls:
  feasible-set size and RTF-minimizing artifact per constraint vector
  (deployment-phase PESQ floor, parameter ceiling, size ceiling).  RTF
  values follow the deployment-phase three-repetition estimand (not
  comparable with Table~\ref{tab:baseline-surface}).}
  \label{tab:joint-selector}
\end{table}

Answering RQ3, Table~\ref{tab:joint-selector} instantiates the selector of
Eq.~\eqref{eq:constrained-selection} over all admitted artifacts of each
lane, using deployment-phase quality, RTF, deployable size, and logical
parameter counts consistently.  Under the moderate operating point
$(q_{\min}{=}3.30, P_{\max}{=}20\mathrm{M}, D_{\max}{=}70~\mathrm{MB})$,
28 B200 artifacts and eight CPU artifacts remain feasible, and the selector
chooses FP16-cast Vocos on B200 and dynamic-INT8 Vocos on CPU; tightening
to $(3.60, 15\mathrm{M}, 55~\mathrm{MB})$ leaves ten and two feasible
artifacts, selecting FP16-cast Vocos on B200 and the untransformed Vocos
control on CPU.  The feasible set remains nonempty on the B200 lane for
quality floors up to RNDVoC's FP16 artifact at 3.894; above that value no
admitted artifact qualifies.
